\section{Related Work}
\label{sec:related-work}

\subsection{Scientific Agent Benchmarks}
Benchmarks for scientific agents have moved from knowledge recall toward
executable research work, and the five sources converted in \benchmark span
the range of task objectives this shift produced. DiscoveryBench frames
data-driven discovery as hypothesis search and verification over provided
datasets, with tasks derived from published findings and graded against
reference hypotheses \citep{majumder2024discoverybench}. BixBench supplies
real bioinformatics analysis capsules and requires agents to execute
multi-step computational biology workflows and answer open-ended questions
about the data \citep{mitchener2025bixbench}. EXP-Bench targets the
experimental process itself, decomposing research experiments extracted
from publications into design and implementation subtasks
\citep{kon2025expbench}. ASI-Bench progressively withdraws methodological
guidance across matched conditions to test autonomous execution
\citep{zhou2026asibench}. SciConvBench, discussed further below, is the
only one of the five whose native format is dialogue rather than artifact
production \citep{somasekharan2026sciconvbench}. Related efforts scale task
collection automatically \citep{li2025autosdt} or test faithful replication
of published analyses \citep{ye2025replicationbench}.

These benchmarks differ substantially in deliverable and grading, which
motivates our source-preserving conversion contract: reducing them to a
common answer format would discard exactly the native scoring that makes a
persona effect interpretable as a scientific error. Their information regimes
also differ: most use one fixed specification, ASI-Bench varies the amount of
methodological guidance, and SciConvBench makes clarification native. None
manipulates collaborator identity and behavior while holding the scientific
task and scoring procedure fixed. \benchmark retains these tasks and their
graders while making that collaborator the experimental variable.

\subsection{Interactive Agents and Simulated Users}
A growing body of work evaluates agents that must interact rather than
execute a complete specification. $\tau$-bench emulates dynamic
conversations between an agent and a language-model user under
domain-specific rules \citep{yao2024taubench}, extended by $\tau^2$-bench to a
dual-control setting in which both agent and user can act on the environment
\citep{barres2025tau2bench}, and AppWorld-UL constructs
tasks that require clarification, confirmation, and infeasibility reporting,
simulating the user with explicitly designed knowledge boundaries
\citep{chen2026appworldul}. In the science and engineering setting,
SciConvBench evaluates multi-turn clarification for ill-posed computational
science requests across fluid mechanics, solid mechanics, materials
science, and PDE formulation, scoring dialogue quality by rubric-based judge and
explicitly requiring no solver, code execution, or tool invocation for scoring
\citep{somasekharan2026sciconvbench}. HiL-Bench isolates the judgment of
\emph{when} to escalate, showing that coding agents which succeed given
complete context degrade sharply when they must themselves decide whether
information is missing \citep{trinh2026hilbench}. Clarification has also
been formalized as a sequential decision problem over whether, what, and
when to ask \citep{ta2026regretbench}, and user simulation has been used to make
agent evaluation more efficient by guiding it toward behavioral diversity
\citep{nakash2026diversity}.

Our work departs from this literature in what is held constant. These
benchmarks fix the human and vary the task, measuring whether the agent
asks well. \benchmark fixes the task and varies the human, measuring how
much of the scientific outcome is attributable to the collaborator. Two
recent lines vary the simulated user deliberately. Non-collaborative user
simulators inject requests for unavailable services, digressions,
impatience, and incomplete utterances, and show state-of-the-art tool
agents degrading substantially under each condition
\citep{shim2025noncollaborative}. Persona Policies argue that LLM-based
simulators inherit the cooperative and homogeneous behavior of their base
models, and search for persona generator programs that induce realistic
behavioral variation while preserving task goals, reporting large fitness
gains and improved agent robustness on $\tau^2$-bench retail and airline
domains \citep{chopra2026ppol}. Together these establish the effect we
investigate: varying the simulated human changes measured agent performance.
Both, however, operate in retail, airline, and customer-service domains,
where a failure is a dialogue breakdown or an incorrect transaction, and both
vary hand-designed or search-discovered behavioral axes rather than
integrated researcher profiles. We instead measure consequences on native
scientific graders, where a collaborator effect surfaces as an incorrect
hypothesis, a flawed analysis, or an unsubmitted deliverable. User simulation for evaluating interactive systems
long predates this literature, from early user models for spoken dialogue
evaluation \citep{eckert1997usermodeling} through agenda-based simulators for
dialogue policy learning \citep{schatzmann2007agenda} to LLM-based behavioral
simulation at scale \citep{wang2023userbehavior}.

\paragraph{Validity of the simulator itself.}
A parallel line of work questions whether LLM user simulators behave as their
specification requires, and it bears directly on our design rather than only on
our citations. Simulators exhibit \emph{assistant bias}, drifting toward helpful
assistant behavior rather than sustaining the user role
\citep{jeong2026assistantbias}---a concern for a researcher simulator whose job
includes gatekeeping and refusal rather than assistance. Bounding what a
simulated user knows is also harder than a prompt instruction suggests, and
naive persona instructions do not reliably constrain a simulator's effective
knowledge \citep{balog2026bounding}. Separately, distributing a task across
multiple turns can reduce model reliability \citep{laban2025lost}, making
interaction length a potential confound for protocols with several review
rounds. We treat these as threats to the manipulation rather than as background,
and \secref{sec:limitations} states how each bounds our claims.

\subsection{Persona-Conditioned Simulation}
Conditioning language models on personal descriptions to elicit distinct
behavior has an established basis. Argyle et al.\ showed that conditioning
on socio-demographic backstories reproduces response distributions of human
subgroups, terming this algorithmic fidelity \citep{argyle2022outofone}. At
scale, Persona Hub demonstrated that a large collection of diverse persona
descriptions can drive correspondingly diverse model outputs
\citep{ge2024personahub}, and subsequent work studies how much generated
diversity persona detail actually produces
\citep{kambhatla2025lexicaldiversity}. For evaluation specifically, Persona
Generators optimize for coverage of the space of possible users rather than
density matching of the most probable ones
\citep{paglieri2026personagenerators}, and frameworks have been proposed to
measure whether persona simulation is faithful along adherence, consistency,
and naturalness \citep{bao2026eval4sim}.

\paragraph{Relation to MatrAIx.}
The researcher profiles used in \benchmark are drawn from MatrAIx, a
population-scale simulated-user evaluation infrastructure
\citep{li2026matraix}. Its Persona 8B corpus represents human variation
through a shared schema of 1,290 categorical dimensions spanning background,
psychology, capability, behavior, and lifestyle. Records are either sampled
from a dependency graph that preserves correlated attributes and enforces
compatibility rules, or extracted from human-authored sources---Wikipedia
biographies, Amazon review histories, the Stack Overflow Developer Survey,
the General Social Survey, PRISM Alignment profiles, and consented survey
responses---mapped into the same schema and de-identified. The upstream release
also attaches natural-language descriptions to human-grounded records, but our
renderer does not use them: the injected identity block is an explicit
attribute--value listing grouped by category, typically on the order of two
hundred populated attributes for a single profile. A condition is therefore a
fully auditable set of stated attributes rather than authored prose, which is
what allows observable behavior to be checked against exact profile text. A
quality-filtered public coreset of
approximately one million records is released, comprising 599,847
human-grounded and 400,000 synthetic personas.

MatrAIx also supplies the strongest existing evidence that this form of
conditioning works: in a 400-trial controlled adherence study spanning ten
behavioral attributes, declared behavior was expressed or correctly
suppressed in 91.5\% of trials \citep{li2026matraix}. That study, however,
was conducted in the four environments MatrAIx provides---Survey, AI
Chatbot, Web, and App---across a task library concentrated in Commerce,
Software, Finance, and Healthcare. Scientific research collaboration is not
among them, and product feedback rather than scientific correctness is the
measured outcome; related demonstrations of persona-based simulation target
the same interactive-application settings \citep{liu2026personaeval}.
\benchmark extends this line of work in exactly that direction. We take the
persona population as given and validated at the level of behavioral
adherence, then ask a question the upstream infrastructure does not: whether
persona-induced behavioral variation propagates into \emph{graded scientific
error} when the simulated human is a research collaborator and the task is
scored by a native scientific verifier. Our contribution is therefore the
scientific environment, protocol, and outcome measurement, not the persona
corpus or its construction.

\paragraph{Caveats adopted from this literature.}
A review of 63 persona studies found that task and target population are
frequently underspecified and that only about a third discuss the
representativeness of their personae, proposing a transparency checklist in
response \citep{batzner2025whosepersonae}. Analytic flexibility in
simulating human data is substantial enough that defensible configuration
choices materially change fidelity conclusions
\citep{cummins2025analyticflexibility}. Work in a scientific domain is
directly instructive: when synthetic personas evaluated a genomics
visualization system, grounded personas shifted toward the concerns of
documented users while ungrounded ones drifted toward issues real experts
never raised, and both missed a preference the expert study observed
\citep{nguyen2026sycamore}. The MatrAIx release is itself explicit that its
human-grounded records are not population-representative, that calibration
targets only four marginal distributions rather than joint coverage over all
1,290 dimensions, and that human studies remain necessary before applying
conclusions to real populations \citep{li2026matraix}. We therefore treat our
profiles as controlled stimuli that induce reproducible interaction
conditions, not as substitutes for human researchers, and we do not interpret
demographic attributes of a profile as properties of any real group.
Additionally, because we inject complete profiles rather than manipulating
individual schema dimensions, communication style, availability, and domain
knowledge vary together within a condition; \secref{sec:limitations} treats
this full-profile confounding directly.

A further distinction from this literature clarifies our architecture.
Persona-conditioned simulation separates the relatively persistent
\emph{persona}---background, identity, values, skills, personality, enduring
preferences---from the more changeable \emph{user state} that many behaviors
actually depend on, including a current goal, available knowledge, prior
interactions, trust, and frustration \citep{liu2026fitnessforpurpose}.
Scientific collaboration requires both: a researcher's disposition is
persistent, but what they know about \emph{this particular study} is not a
property of their biography. \benchmark separates these sources accordingly.
The persona supplies persistent disposition and is the manipulated variable;
study-specific context is owned by the task and held identical across
persona conditions; and private evaluator references are withheld from the
researcher entirely. A profile therefore changes how study context is
communicated and reviewed, not what ground truth exists, which is what makes
a score difference interpretable as a consequence of collaboration rather
than of unequal information endowment.

This separation is imperfect in one respect worth stating. The MatrAIx schema
is not purely dispositional: rendered profiles include a small
\emph{interaction state} group with fields such as the participant's expertise
gap relative to the system, work schedule, and work location. At least one such
field is potentially consequential, since a profile that declares its holder an
expert testing the system frames the interaction differently from one seeking
genuine assistance. These fields are part of the injected condition and are not
separately manipulated, so the persona conditions carry a limited amount of
state alongside disposition.

\subsection{Evaluating Collaboration}
Evaluations of interactive agents tend to report either outcome scores or
interaction quality, and the two can diverge. Outcome-only scoring cannot
distinguish an agent that obtained the information it needed from one that
was simply given an easier collaborator; interaction-only scoring, as in
rubric-graded clarification dialogues, establishes that an agent asked
sensible questions without establishing that the resulting science was
correct \citep{somasekharan2026sciconvbench}. Metrics that combine the two,
such as HiL-Bench's pairing of task success with the precision and recall
of escalation decisions, come closer by asking whether the agent asked
about the things that actually mattered \citep{trinh2026hilbench}.

\benchmark reports native outcome scores and interaction metrics as separate
quantities over the same episodes, with a paired persona-off baseline and
repeated independent episodes per cell. This supports three claims that a
single aggregate cannot: that a collaborator changed observable agent
behavior, that a collaborator changed graded scientific correctness, and
that these two need not co-occur. We additionally analyze stage-aligned
conversation trajectories and present evidence-linked cases tracing profile
wording through researcher behavior to actor decisions and graded outcomes.

\paragraph{Which validity we claim.}
A fitness-for-purpose framework for persona-conditioned simulation argues
that simulated users should be judged against the specific claim they were
built to support rather than against general human-likeness, and distinguishes
four forms of validity: persona fidelity, whether the model preserves and
uses the supplied profile; behavioral and interaction validity, whether
actions and trajectories resemble the target human process; task and
ecological validity, whether evaluation with simulated users exposes the same
capabilities and failures real users would reveal; and distributional
validity, whether a collection of simulated users represents an intended
population \citep{liu2026fitnessforpurpose}. Passing one form does not imply
the others, and the framework's central recommendation is to state exactly
what has been validated \citep{li2026humansimulation}. We adopt that
discipline. \benchmark inherits persona fidelity from upstream adherence
validation \citep{li2026matraix} rather than re-establishing it, and its own
contribution is evidence about task validity in a new domain: whether
collaborator variation exposes scientific failures that a single default
collaborator does not. We explicitly do not claim distributional validity.
The final profile panel was selected using pilot outcomes, so it characterizes
the range of collaborator effects rather than estimating an average effect
over any researcher population, and we make no claim that behavioral and
interaction validity holds against real scientist behavior, which would
require human comparison data we do not have. \secref{sec:limitations}
develops these constraints.
